\section{\heiti 实验}
\label{sec:experiments}

本节在 ImageNet-100 上做受控实验。基准配置取 240M 档：编码器与解码器都是标准 Transformer（softmax 注意力），可训参数分别为 244.43 M 与 237.11 M，体量基本对齐，因此不同设置之间的差异可以归因到结构因素；除特别说明外，所有比较都在同一训练协议内进行。这一档位不是预先指定的，而是由第~\ref{sec:ksweep}~节的容量阶梯选出——在本协议下容量收益在 240M 档达到峰值，继续加宽反而回落。

\subsection{\heiti 实验设置}
\label{sec:impl}

\textbf{数据与模型.} 数据为 ImageNet-100，训练集 121,689 张、验证集 5,000 张；输入 $224\times224$，patch 为 $16\times16$，故 $N=196$；token 预算 $K=32$，以 token 数计的压缩比 $N/K=6.12$。编码器为 13 层 Transformer，隐层维度 1248、16 头，可训参数 244.43 M；解码器为 7 层 Transformer 解码器堆叠，隐层维度 1440、16 头，可训参数 237.11 M；两者合计 481.54 M，体量相当（244.43 M 对 237.11 M），因此两者之间的比较不涉及“模型更大”的解释。每个特殊 token 是 1248 维连续向量，32 个 token 共 39{,}936 个连续值。主要配置见表~\ref{tab:config}。

\textbf{训练与指标.} 训练分两阶段。阶段 1 为单步解码：$T=1$、读窗口取 $[0,K]$（一次读完全部 32 个特殊 token），45 个 epoch、batch size 128、共 42{,}750 步，从零训练；阶段 2 为 16 步循环：从阶段 1 的最优权重热启动，采用均匀窗口（每步 2 个特殊 token），6 个 epoch、batch size 64、共 11{,}406 步，编码器与解码器联合微调。优化器为 AdamW，学习率 $2\times10^{-4}$，cosine 调度并带 8\% warmup，权重衰减 0.05，梯度裁剪 1.0，全程 bf16。评价指标以峰值信噪比（PSNR）为主：在 $[0,1]$ 空间中逐图计算、再对 5,000 张留出图取平均，单位 dB，越高越好；辅以结构相似度（SSIM）与像素 L1（0--255 空间，不裁剪）。三者都是逐图配对指标，直接比较每张重建图与它对应的 GT，衡量的是信息被保留下来的程度。本文不以重建 FID（rFID）为主指标：它以 InceptionV3 特征在整批重建图与整批 GT 图之间比较均值与协方差，是 unpaired 的集合级分布度量，不含逐图重建精度（一批“逐图全错、但都取自自然图像分布”的输出也能得到很低的 rFID），且对纹理统计的权重远大于几何误差；它适合评价生成式模型的真实感，不适合本文的信息保持问题。

\textbf{瓶颈臂.} 为检验每 token 维数的影响，本文在 240M 档上另训两支瓶颈臂：在编码器末端加一层线性映射，把 1248 维的特殊 token 分别压到 448 维与 224 维（$z_s$ 的存储值数由 39,936 降到 14,336 与 7,168，即 2.79 倍与 5.57 倍）；两支臂均从 16 步模型热启动、全量微调 3 个 epoch（batch size 64、学习率 $2\times10^{-4}$、cosine 调度，其余同阶段 2），可训参数分别为 481.51 M 与 480.63 M，与主模型（481.54 M）基本一致。每步仍读 2 个特殊 token，故第 $t$ 步累计 $896t$ 与 $448t$ 个连续值，16 步的上限分别为 14{,}336 与 7{,}168 个值。

\textbf{参考基线.} 本文不使用量化与熵编码，“压缩率”以 token 数而非比特数度量。不随模型结构变化的两条像素基线为：逐 patch 直流（每个 patch 用其均值表示，18.45 dB）与常量灰（11.81 dB），用于判断模型是否真正学到了图像信息。第~\ref{sec:compare}~节另用逐块 PCA 作为有学习基的线性对照：基在训练集 patch（400 张图的 78{,}400 个 patch）上学习、全图共享、不逐图拟合，每块保留前 $r$ 个主成分，故每图信息量为 $196r$ 个连续值。

\textbf{口径说明.} 本文主实验的 16 个预算依次为 $2,4,\dots,32$ 个特殊 token（第 $t$ 步累计 $2t$ 个）。对照统一按\textbf{每图连续值个数}对齐：主模型每 token 1248 维，32 个 token 共 39{,}936 个值；448/224 维瓶颈臂每图最多 14{,}336 与 7{,}168 个值（第 $t$ 步累计 $896t$ 与 $448t$）；逐块 PCA 每图 $196r$ 个值，取 $r=9/18/37/73$ 与瓶颈臂对齐（1{,}764/3{,}528/7{,}252/14{,}308，相对差 $\le1.6\%$）。图~\ref{fig:complexity}、图~\ref{fig:recon} 取自主模型（$K=32$、$T=16$），图~\ref{fig:step} 为 224 维瓶颈臂。压缩率以 token 数度量：旧 63M 基准为 $3.06\times$（196$\to$64），本文 240M 档为 $6.12\times$（196$\to$32）；以连续值个数计时，主模型每图 39{,}936 个值、瓶颈臂为 14{,}336/7{,}168 个值——两种口径在文中分别标明，不混用。

\begin{table}[htbp]
\centering\footnotesize\setlength{\tabcolsep}{3pt}
\caption{实验配置}
\label{tab:config}
\begin{tabular}{ll}
\toprule
项目 & 设置 \\
\midrule
数据 & ImageNet-100（121,689 / 5,000） \\
输入 / patch & $224\times224$ / $16\times16$，$N=196$ \\
token 预算 & $K=32$，压缩比 $N/K=6.12$ \\
每图信息量 & $32\times1248=39{,}936$ 个连续值 \\
编码器 & 13 层，$D=1248$，16 头，244.43 M \\
解码器 & 7 层，$D'=1440$，16 头，237.11 M \\
总参数 & 481.54 M \\
阶段 1 & 单步（$T=1$），45 epoch，bs 128，42,750 步 \\
阶段 2 & 16 步（每步 2 token），6 epoch，bs 64，11,406 步 \\
优化器 & AdamW，lr $2\times10^{-4}$，warmup 8\% \\
精度 / 硬件 & bf16 / A800 80 GB \\
评价指标 & PSNR（主）、SSIM、L1（0--255） \\
\bottomrule
\end{tabular}
\end{table}

\subsection{\heiti 与线性码的预算分层对比}
\label{sec:compare}

本节按“每图连续值个数”把本文的码与逐块 PCA 基线对齐：同一份 5,000 张留出图、同一评测口径（逐图 PSNR 为主，SSIM 与像素 L1 为辅），数值见表~\ref{tab:budget} 与图~\ref{fig:step}。为覆盖不同预算区间，这里同时使用三支模型：主模型（每 token 1248 维，最多 39{,}936 个值）与两支瓶颈臂（每 token 448/224 维，最多 14{,}336/7{,}168 个值，见第~\ref{sec:impl}~节）。其他分词器不在本表并列：TiTok 与 ALIT 都带 VQGAN 量化底座，ALIT 还含熵编码与对抗损失，训练数据、重建分辨率与留出集也各不相同，数值直接并列不可靠；本文不作跨方法的数值对比，待按统一协议复现其训练与评测后再行比较。

\begin{table}[!t]
\centering\scriptsize\setlength{\tabcolsep}{3pt}
\caption{预算分层的同值数 PSNR 对照（dB）}
\label{tab:budget}
\begin{tabular}{lcccc}
\toprule
\multirow{2}{*}{方法} & \multicolumn{4}{c}{每图连续值个数} \\
\cmidrule(lr){2-5}
 & 1{,}792 & 3{,}584 & 7{,}168 & 14{,}336 \\
\midrule
本文 448 维臂（4/8/16/32 token） & 23.19 & 24.99 & 26.63 & 27.69 \\
本文 224 维臂（8/16/32 token） & 23.66 & 25.21 & 26.08 & / \\
逐块 PCA（rank 9/18/37/73） & 22.26 & 24.10 & 26.28 & 29.06 \\
\bottomrule
\end{tabular}

\vspace{0.5mm}
{\scriptsize
注：5,000 张留出图上逐图 PSNR 的均值，越高越好。两支瓶颈臂在编码器末端各加一层线性映射，把 1248 维的特殊 token 压到 448 / 224 维；每步仍读 2 个 token，故第 $t$ 步累计 $896t$ / $448t$ 个值。列头四档分别对应 448 臂的 $t=2,4,8,16$ 与 224 臂的 $t=4,8,16$（224 臂的上限为 7{,}168 值）。逐块 PCA 每图为 $196r$ 个值，$r=9/18/37/73$ 对齐四档（1{,}764/3{,}528/7{,}252/14{,}308，相对差 $\le1.6\%$），共享基（$768\times256$）算模型参数、不计入每图预算。相同值数下两臂读入的 token 数不同（7{,}168 值处 448 臂 16 个、224 臂 32 个），比较时须注意这一结构差异。}
\end{table}

\begin{figure*}[t]
\centering
\includegraphics[width=0.95\textwidth]{图/fig_step_pca_zs224.pdf}
\caption{按读出步数比较像素 L1（左）与 PSNR（右）：本文 224 维臂与同值数逐块 PCA。横轴为读出步数，每步读 2 个 224 维 token（448 个值），第 $t$ 步累计 $448t$ 个值；PCA 按 $r=448t/196$ 在实测的 rank 1--40 之间线性插值，与本文严格同值数。实线为本文，虚线为逐块 PCA；L1 上本文 16 步全程更低，PSNR 上第 1--15 步本文更高、约在第 15.1 步交叉。}
\label{fig:step}
\end{figure*}

由表~\ref{tab:budget} 与图~\ref{fig:step} 可得四点。第一，\textbf{紧预算端本文的码更好}：在 1{,}792--3{,}584 个连续值处，448 维臂为 23.19 / 24.99 dB、224 维臂为 23.66 / 25.21 dB，而同值数 PCA（rank 9/18）为 22.26 / 24.10 dB，即本文高 $+0.89\sim+1.41$ dB；L1 同向（11.78 / 11.24 对 13.79 px，9.67 / 9.52 对 11.21 px）。原因是预算极紧时逐块 PCA 的每个 patch 只保留 9--18 个主成分、几乎没有容量，而本文的码把整幅图像的信息分摊到 32 个 token 上，并由像素监督学到跨位置的先验，因此在极紧预算下更省。

第二，\textbf{预算放宽后线性码反超}：7{,}168 值处 448 维臂仍高 0.34 dB（26.63 对 26.28 dB），224 维臂的 PSNR 为 26.08 dB、低 0.20 dB，但其 L1 低 0.18 px（8.71 对 8.89 px）；到 14{,}336 值处 PCA 反超 1.37 dB（29.06 对 27.69 dB）。主模型自身在同值数下也不占优：在 4{,}992--39{,}936 值的四档上 PCA 均更高，39{,}936 值处为 27.94 dB 对 35.75 dB。这与“逐块 PCA 保留 196 个空间位置的码、只降低每个位置的维数”一致：预算越松，它越接近逐位置线性编码的优势区间。

第三，图~\ref{fig:step} 把 224 维臂与同值数 PCA 按读出步数逐点比较：每步读 2 个 224 维 token（448 个值），第 $t$ 步累计 $448t$ 个值；PCA 在实测的 rank 1--40 之间按 $r=448t/196$ 线性插值，与本文严格同值数。L1 上本文 16 步全程更低，优势由第 1 步的 $-6.84$ px 单调收窄到第 16 步（7{,}168 值）的 $-0.22$ px，始终没有交叉；PSNR 上第 1--15 步本文更高（$+2.63$ 降到 $+0.02$ dB），第 16 步 PCA 反超 0.17 dB，交叉点约在第 15.1 步。两个指标的交叉位置不同，是因为 PSNR 以均方误差为底、对少数大误差像素更敏感，而 L1 是逐像素平均；因此“紧预算端本文更好”的成立范围与指标有关：按 L1 可写到 7{,}168 值，按 PSNR 只到约 6{,}720 值（第 15 步）。

第四，\textbf{压窄每 token 的维数代价很小}：把每 token 由 1248 维压到 448 维后，每图连续值由 39{,}936 降到 14{,}336（36\%），末步 PSNR 仅由 27.94 降到 27.69 dB；在同一 14{,}336 值下比旧 63M 基准（24.15 dB）高 3.54 dB。这说明 1248 维中的冗余较大、读出质量的主要瓶颈在 token 数而非每 token 的宽度。相同值数下“少而宽”与“多而窄”的取舍并不对称：7{,}168 值处 448 维臂（16 个 token）为 26.63 dB，高于 224 维臂（32 个 token）的 26.08 dB，说明在本协议里把预算摊到更多 token 上并不必然更好。

因此本文的定位不是“同值数下最优的像素重建器”，而是“单一模型中可截断的 32-token 语义码”：在不增加参数、不重新训练的前提下，一个 32-token 的表示覆盖由少到多的多档预算，而逐块 PCA 既没有减少 token 数，也不能按步截断到更少预算；在极紧预算端（$\le3{,}584$ 值）本文的码更好，在较松预算端线性码更好。本文暂不与其他分词器并列数值：TiTok、ALIT 等带量化码本、熵编码与对抗损失，训练数据、分辨率与评测集都不同，直接比较不可靠；待按统一协议复现其训练与评测后再行比较。

\subsection{\heiti 图像复杂度}
\label{sec:complexity}

读出步数增加会改善重建，但不同图像是否同等受益、容量又如何影响这一差异，仍需检验。本小节使用公开的 SAVOIAS 人评复杂度数据\cite{saraee2018savoias}作为分布外分析集：People-Art 420 张、Places 200 张、MSCOCO 200 张，复杂度分数为 0--100 的众包成对比较结果。被测模型为 ImageNet-100 上训练的 240M 档（$K=32$、$T=16$），未在 SAVOIAS 上微调。为把“容量”与“预算”两个因素分开，本小节同时用到第~\ref{sec:ksweep}~节的 5 个容量档（合计 63.0 M、119.3 M、238.2 M、481.5 M、720.6 M），它们与基准同协议（$K=32$、$T=16$、每步 2 个 token），唯一自变量是宽度。

\begin{figure*}[t]
\centering
\includegraphics[width=0.95\textwidth]{图/fig_savoias_heatmap_m240.png}
\caption{复杂度热图（240M 档，16 档预算）}
\label{fig:complexity}
\end{figure*}

图~\ref{fig:complexity} 给出 240M 档的热图：列为 People-Art、Places、MSCOCO，行为 token 预算（2 到 32，第 $t$ 步累计 $2t$ 个特殊 token），横轴为人评复杂度（15 到 85），颜色为像素 L1（0--255，逐图高斯平滑，$\sigma=5$），越深误差越大。两个方向与旧协议（63M/K=64）一致：固定 token 预算时，复杂度越高误差越大；固定复杂度时，token 越多误差越低，且高复杂度一列下降得更快。

三个子集上复杂度与重建误差在全部 16 个预算上正相关：最少预算（2 token）时 $\rho=+0.520$、$+0.576$、$+0.415$（People-Art、Places、MSCOCO），最多预算（32 token）时 $\rho=+0.518$、$+0.477$、$+0.327$；$\rho$ 随预算只轻微下降，说明复杂度效应不是某一档预算的偶然。把容量维度也纳入（第~\ref{sec:ksweep}~节的 5 个容量档 × 3 个子集 = 15 个组合），15 格全部为正（$\rho\in[+0.327,+0.536]$），并随容量只有轻微下降（子集平均从 63M 档的 0.489 降到 482M 档的 0.441，721M 档回升到 0.453）。

按复杂度三分位分组，240M 档上最高组与最低组的末步 L1 差为 3.629 px（People-Art）、2.821 px（Places）、1.504 px（MSCOCO）。容量增加缩小的是这一差的绝对幅度：从 63M 档的 5.655 / 4.992 / 2.811 px 降到 482M 档的 3.629 / 2.821 / 1.504 px（降幅 36\% / 43\% / 47\%）；但相对幅度基本不变（482M 档分别为 66.7\%、49.7\%、22.5\%）。也就是说，容量买到的是整体误差水平的下降，“越复杂越难”的相对结构没有被容量抹平。

在分布外数据上，240M 档末步的 L1 为 7.231、PSNR 为 27.406 dB、SSIM 为 0.8257，与同协议 ImageNet-100 验证集上的 7.004、27.936 dB、0.8288 接近（L1 差 $+0.227$ px、PSNR 差 $-0.53$ dB）；OOD 上的最优停步为第 15 步，与 ID 一致（5 个容量档依次为 11 / 13 / 13 / 15 / 16 步）。逐步曲线在第 15 步之后基本饱和，说明该收敛位置由读出协议（每步 2 个 token、$T=16$）而非数据分布决定。

图~\ref{fig:recon} 给出低、中、高复杂度图像在 240M 档上的重建示例：三行的复杂度依次为 10、48、89（文件 138/15/21.jpg），列为原始图像与第 1、2、4、8、16 步的重建（累计 2、4、8、16、32 个 token）。复杂度 10 的图像在第 1 步（2 token）已还原主要物体与布局（L1 5.27 px），继续读入主要改善边缘与纹理（末步 2.38 px）；复杂度 48 与 89 的图像在第 1 步只有色块与大致明暗（14.53 / 14.81 px），到第 4 步（8 token）区域边界与物体轮廓才清晰（9.76 / 10.69 px），末步降到 6.50 / 7.60 px。低复杂度行的相对降幅与另两行相近（55\% 对 55\%、49\%），但绝对改善明显更小（2.89 对 8.03、7.21 px）：复杂度主要抬高误差的起点与整体水平，读出步数的收敛节奏没有被改变。

\begin{figure*}[t]
\centering
\includegraphics[width=0.95\textwidth]{图/fig_recon_grid_m240.pdf}
\caption{重建示例（240M 档，$K=32$、$T=16$；列为原始图像与第 1、2、4、8、16 步（2--32 token）的重建）}
\label{fig:recon}
\end{figure*}

\subsection{\heiti 容量阶梯}
\label{sec:ksweep}

本文的基准配置由容量阶梯选出。前面的实验固定 $K=32$、$T=16$、每步 2 个特殊 token，本小节只改变编码器与解码器的宽度，得到 5 个容量档：编码器与解码器各约 30 M（合计 63.0 M）、60 M（119.3 M）、120 M（238.2 M）、240 M（481.5 M）与 360 M（720.6 M）。深度固定（编码器 13 层、解码器 7 层），只加宽隐层；5 档共用同一训练日程（阶段 1 单步 45 epoch、阶段 2 16 步 6 epoch），容量是唯一自变量。结果见图~\ref{fig:scale}。

\begin{figure}[t]
\centering
\pgs{\begin{tikzpicture}
\begin{axis}[
  name=capa, width=0.78\linewidth, height=5.5cm, scale only axis,
  xlabel={总参数 / M}, ylabel={PSNR / dB},
  symbolic x coords={63,119,238,482,721}, xtick=data,
  ymin=23.8, ymax=28.4,
  tick label style={font=\scriptsize}, label style={font=\footnotesize},
  axis y line*=left,
  legend style={font=\scriptsize,draw=none,fill=none,at={(0.02,0.98)},anchor=north west},
  legend cell align=left,
]
\addplot[color=blue!60!black,line width=0.9pt,mark=*,mark size=1.6pt]
  coordinates {(63,24.15) (119,25.75) (238,27.31) (482,27.94) (721,27.31)};
\addlegendentry{PSNR（左轴）}
\addlegendimage{color=orange!85!black,line width=0.9pt,dashed,mark=square*,mark size=1.4pt}
\addlegendentry{SSIM（右轴）}
\end{axis}
\begin{axis}[
  at={(capa.south west)}, anchor=south west,
  width=0.78\linewidth, height=5.5cm, scale only axis,
  symbolic x coords={63,119,238,482,721}, xtick=\empty,
  ymin=0.64, ymax=0.845,
  axis y line*=right, axis x line*=none, yticklabel pos=right,
  tick label style={font=\scriptsize}, label style={font=\footnotesize},
]
\addplot[color=orange!85!black,line width=0.9pt,dashed,mark=square*,mark size=1.4pt]
  coordinates {(63,0.664) (119,0.739) (238,0.804) (482,0.829) (721,0.802)};
\end{axis}
\end{tikzpicture}}
\caption{容量阶梯}
\label{fig:scale}
\end{figure}

$T=16$ 末步的 PSNR 依次为 24.148、25.747、27.308、27.936、27.314 dB，归一化空间的 L1 依次为 0.1877、0.1570、0.1320、0.1215、0.1305，SSIM 依次为 0.6638、0.7389、0.8039、0.8288、0.8015；相邻档的边际收益为 $+1.60$、$+1.56$、$+0.63$、$-0.62$ dB，按每 100 M 参数折算依次为 $+2.84$、$+1.31$、$+0.26$、$-0.26$ dB/100M。由此得到三点。第一，\textbf{容量甜点在 240M 档（合计 481.54 M）}：从 63M 到 240M 用 418 M 参数买来 $+3.79$ dB，而从 240M 到 360M 用 239 M 参数反而亏 0.62 dB；L1 与 SSIM 同向，说明这不是单一指标的噪声。第二，\textbf{阶段 1 与阶段 2 同序}：两阶段的末步 PSNR 分别为 28.728 / 30.242 / 31.205 / 32.345 / 31.977 dB 与 24.148 / 25.747 / 27.308 / 27.936 / 27.314 dB，B 相相对 A 相的落差在 4.10--4.66 dB 之间几乎不随容量变化，说明“每步 2 个 token、16 步”的窄窗口读出损失不是容量问题，加参数填不平。第三，\textbf{容量抬高的是每一步的整体水平，而非收敛速度}：第 1 步的 L1 从 17.18 px 降到 13.50 px，而最优停步只在 11 / 13 / 13 / 15 / 16 之间后移。

作为旁证，在更早的 63M/K=64、8 步递增窗口协议下还做过 token 预算的扫描（$K=8$ 到 128，共 9 档，每档单独训练）：$K\le64$ 时质量单调上升（PSNR 由 21.78 dB 升到 27.76 dB），$K>64$ 后回落（$K=96$ 为 26.24 dB、$K=128$ 为 25.09 dB），因此旧协议的预算操作点是 $K=64$。该结果与新基准不可直接相减（模型宽度、步数与窗口计划都不同），只作为“预算需与容量、步数匹配”的历史旁证保留；新基准下预算与步数的匹配关系是 $K=2T$。
